\subsection{COXETER TRANSFORMATION}

Let C be a chamber of R, let \(\{\alpha_{1},\ldots,\alpha_{l}\}\) be the corresponding basis of R, and let \(c=s_{\alpha_{1}}\ldots s_{\alpha_{l}}\) . The element c of W is called the Coxeter transformation of W defined by C and the bijection \(i\mapsto\alpha_{i}\) (Chap. V, § 6, no. 1). Its order h is called the Coxeter number of W (or of R).

PROPOSITION 30. Assume that R is irreducible. Let m be an integer between 1 and h - 1 and prime to h. Then \(\exp\left(\frac{2i\pi m}{h}\right)\) is an eigenvalue of c of multiplicity 1.

In particular, \(m\) is an exponent of \(W\) , cf.~Chap. V, § 6, no. 2.

We first prove a lemma:

Lemma 4. For all \(w \in W\) , the characteristic polynomial of \(w\) has integer coefficients.

We know (no. 6, Th. 3) that \(\{\alpha_{1},\ldots,\alpha_{l}\}\) is a basis of the subgroup \(Q(R)\) of V generated by R. Since w leaves \(Q(R)\) stable, its matrix with respect to \(\{\alpha_{1},\ldots,\alpha_{l}\}\) has integer entries; hence its characteristic polynomial has integer coefficients.

Let P be the characteristic polynomial of c.~The above lemma shows that the coefficients of P are integers. By Chap. V, § 6, no. 2, Cor. 2 of Prop. 3, the primitive hth root of unity \(z = \exp\left(\frac{2i\pi}{h}\right)\) is a simple root of P. Every conjugate of z over Q is therefore also a simple root of P. But we know (Algebra, Chap. V) that the primitive hth roots of unity are conjugate over Q. They are thus all simple roots of P, which proves the proposition.

PROPOSITION 31. Assume that R is irreducible and reduced, and let \(\beta = n_{1}\alpha_{1} + \cdots + n_{l}\alpha_{l}\) be the highest root of R (cf.~no. 8). Then \(n_{1} + \cdots + n_{l} = h - 1\) .

Let \(R_{+}\) be the set of positive roots relative to C. Then (no. 10, Cor. of Prop. 29):

\[
\begin{array}{l} n _ {1} + \dots + n _ {l} = \frac {1}{2} \sum_ {\alpha \in \mathrm{R} _ {+}} n (\beta , \alpha) \\ = 1 + \frac {1}{2} \sum_ {\alpha \in \mathrm{R} _ {+}, \alpha \neq \beta} n (\beta , \alpha) = 1 + \sum_ {\alpha \in \mathrm{R} _ {+}, \alpha \neq \beta} \frac {(\alpha | \beta)}{(\alpha | \alpha)}. \end{array}
\]

By no. 8, Prop. 25 (iv), for all \(\alpha \in \mathbb{R}_+\) and \(\alpha \neq \beta\) , \(n(\alpha, \beta) = 0\) or 1, so \(n(\alpha, \beta)^2 = n(\alpha, \beta)\) , that is \(\frac{4(\alpha | \beta)^2}{(\beta | \beta)^2} = \frac{2(\alpha | \beta)}{(\beta | \beta)}\) . Hence:

\[
\begin{array}{l} n _ {1} + \dots + n _ {l} + 1 = 2 + 2 \sum_ {\alpha \in R _ {+}, \alpha \neq \beta} \frac {(\alpha | \beta) ^ {2}}{(\alpha | \alpha) (\beta | \beta)} \\ = 2 \sum_ {\alpha \in R _ {+}} \frac {(\alpha | \beta) ^ {2}}{(\alpha | \alpha) (\beta | \beta)} = (\beta | \beta) ^ {- 1} \sum_ {\alpha \in R} (\frac {\alpha}{\| \alpha \|} | \beta) ^ {2}. \end{array}
\]

By Chap. V, § 6, no. 2, Cor. of Th. 1,

\[
\sum_ {\alpha \in R} \left(\frac {\alpha}{\| \alpha \|} | \beta\right) ^ {2} = h (\beta | \beta)
\]

\[
\mathrm{so} n _ {1} + \dots + n _ {l} + 1 = h.
\]

PROPOSITION 32. Assume that R is irreducible, and that all roots have the same length. Let \(\alpha \in R\) . The number of elements of R not orthogonal to \(\alpha\) is 4h - 6.

Let \(R'\) be the set of roots not proportional to and not orthogonal to \(\alpha\) . By Chap. V, § 6, no. 2, Cor. of Th. 1,

\[
(\alpha | \alpha) ^ {2} + (\alpha | - \alpha) ^ {2} + \sum_ {\beta \in \mathrm{R} ^ {\prime}} (\alpha | \beta) ^ {2} = h (\alpha | \alpha) ^ {2},
\]

that is

\[
\sum_ {\beta \in \mathsf {R} ^ {\prime}} (\alpha \mid \beta) ^ {2} = (h - 2) (\alpha \mid \alpha) ^ {2}.
\]

If \(\beta \in \mathbf{R}'\) , then \((\alpha | \beta) = \pm \frac{1}{2} (\alpha | \alpha)\) by the list in no. 3. Hence

\[
\frac {1}{4} \operatorname{Card} R ^ {\prime} = h - 2, \quad \operatorname{Card} R ^ {\prime} = 4 h - 8,
\]

and the number of roots not orthogonal to \(\alpha\) is \(\operatorname{Card} R' + 2 = 4h - 6\) .

PROPOSITION 33. Assume that R is irreducible and reduced. Put \(s_{\alpha_i} = s_i\) , and let \(\Gamma\) be the subgroup of W generated by \(c = s_1 \ldots s_l\) .

\begin{enumerate}
\def\labelenumi{(\roman{enumi})}
\item
  Let \(\theta_{i} = s_{l}s_{l - 1}\dots s_{i + 1}(\alpha_{i})\) ( \(i = 1,\dots ,l\) ). Then, \(\theta_{i} > 0,c(\theta_{i}) < 0\) .
\item
  If \(\alpha\) is a root \(>0\) such that \(c(\alpha) < 0\) , then \(\alpha\) is equal to one of the \(\theta_{i}\) .
\item
  The family \((\theta_{i})_{1\leqslant i\leqslant l}\) is a basis of \(\mathrm{Q}(\mathrm{R})\)
\item
  Let \(\Omega_{i}\) be the orbit of \(\theta_{i}\) under \(\Gamma\) . The sets \(\Omega_{i}\) are pairwise disjoint, they are all the orbits of \(\Gamma\) on \(\mathbb{R}\) , and each has \(h\) elements.
\end{enumerate}

Observe first that \((s_1, \ldots, s_l)\) is a reduced decomposition of \(c\) (Chap. IV, § 1, no. 1) with respect to the set \(S\) of the \(s_i\) . Indeed, otherwise there would exist a subset \(X = S - \{j\}\) of \(l - 1\) elements of \(S\) such that \(c \in W_X\) , which would contradict Cor. 2 of Prop. 7 of Chap. IV, § 1, no. 8.

Applying Cor. 2 of Prop. 17 of no. 6 to \(c\) gives assertions (i) and (ii).

Let \(Q_{i}\) be the subgroup of \(Q(R)\) generated by the \(\alpha_{j}, j > i\) . It is immediate that \(Q_{i}\) is stable under the \(s_{j}, j > i\) , and that \(s_{j}(\alpha_{i}) = \alpha_{i} \mod Q_{i}\) for \(j > i\) . Hence:

\[
\theta_ {i} = s _ {l} \dots s _ {i + 1} (\alpha_ {i}) = \alpha_ {i} \bmod . Q _ {i}.
\]

In other words, there exist integers \(c_{ij}\) such that

\[
\theta_ {i} = \alpha_ {i} + \sum_ {j > i} c _ {i j} \alpha_ {j}.
\]

Part (iii) follows immediately.

Finally, let \(\alpha\) be a root. The element \(\sum_{k=0}^{h-1} c^k(\alpha)\) is invariant under \(c\) , and hence zero (Chap. V, § 6, no. 2). The \(c^k(\alpha)\) cannot therefore all have the same sign, and there exists a \(k\) such that \(c^k(\alpha) > 0\) and \(c^{k+1}(\alpha) < 0\) . By (ii), \(c^k(\alpha)\) is one of the \(\theta_i\) . Thus every orbit of \(\Gamma\) on R is one of the \(\Omega_i\) . Extend \((x|y)\) to a hermitian form on V \(\otimes\) C. Each orbit of \(\Gamma\) in R has at most \(h\) elements, and there are at most \(l\) distinct orbits, by the above remarks. Now (Chap. V, § 6, no. 2, Theorem 2, ii)), the cardinal of R is equal to \(hl\) , which immediately implies (iv).
% semantic-fragments: legacy-input-seam
\relax{}
